\section{Related Work}

\subsection{Real-Time and Small-Object Detection}
The COCO evaluation protocol established averaged AP over intersection-over-union (IoU) thresholds as a standard measure of localization quality \cite{lin2014coco}. For small objects, however, repeated downsampling destroys both energy and phase information before a detector can form a stable response. High-resolution pyramids mitigate this loss, but their cost scales with spatial area. YOLOv8 provides a compact anchor-free implementation and reproducible training stack \cite{ultralytics2023}. Because SRPA-YOLO is developed within this framework, YOLOv8n is the primary structural and deployment baseline. YOLO11 is included as a contemporary compact control, whereas UAV-specific reproductions test whether the proposed design remains competitive beyond its parent architecture.

Hardware-aware modules reduce the cost of retaining high-resolution features. MobileNet replaces dense spatial convolution with depthwise separable operations \cite{howard2017mobilenets}; GhostNet synthesizes redundant features through cheap transforms \cite{han2020ghostnet}; FasterNet argues that realized FLOPS and memory access are as important as nominal FLOPs \cite{chen2023fasternet}. Our backbone and neck use C2f aggregation, Ghost convolution, depthwise downsampling, and a partial-convolution stage in this spirit. CBAM combines channel and spatial selection \cite{woo2018cbam}, while efficient multi-scale attention (EMA) captures cross-spatial relationships with limited parameters \cite{ouyang2023ema}. These modules are retained only where the measured architecture benefits; the paper's novelty is not an attention replacement.

\subsection{Anti-UAV Benchmarks and Specialized Detectors}
Anti-UAV benchmarks have progressively expanded visible and thermal tracking conditions \cite{zhao2022antiuav,jiang2023antiuav,huang2024antiuav410}. The DUT Anti-UAV detection subset is especially suitable for evaluating long-range single-class UAV detection because it contains large background regions and extreme target-scale variation. Specialized detectors have addressed these conditions through stronger feature fusion, lightweight blocks, and attention. YOLOv7-GS modifies the real-time detector for cluttered anti-UAV scenes \cite{bo2024yolov7gs}; EDGS-YOLOv8 and DRBD-YOLOv8 combine lightweight operations with enhanced multiscale representation \cite{huang2024edgs,jiang2024drbd}; RLRD-YOLO emphasizes the balance between a lightweight model and robust small-object recognition \cite{li2025rlrd}.

IASL-YOLO combines a lightweight backbone with adaptive feature fusion, whereas BRA-YOLOv10 couples multiscale representation with structured channel pruning \cite{yang2025iasl,zhang2025brayolo}. Both methods improve UAV detection by reallocating computation in the backbone or feature pyramid. SRPA-YOLO addresses a complementary source of inefficiency inside the prediction head: it constrains the classification and localization feature row spaces, assigns a low-rank classification-private representation, and removes the training-time decomposition through exact fusion. The distinction is therefore architectural---feature-extraction and pyramid efficiency in the former methods versus task-specific head-capacity allocation and static deployment equivalence in SRPA-YOLO.

Broader aerial benchmarks such as VisDrone contain multiple object classes and denser scenes \cite{zhu2018visdrone}. Their task semantics differ from the single-class long-range UAV detection problem considered here, so published results on those benchmarks are not directly comparable with the present DUT Anti-UAV evaluation. Evaluation under additional target distributions remains a direction for future work.

\subsection{Task Conflict, Distillation, and Deployable Training}
Dense heads often share features before separate output layers. TSD and TOOD show that this sharing can create misaligned optima between class confidence and box localization \cite{wu2020tsd,feng2021tood}. Our response is not a fully decoupled head; it is a low-rank task allocation in which expensive spatial representation remains mostly shared. The classification-private representation is additive and classification-only, so the box Jacobian with respect to this representation is identically zero.

Knowledge distillation transfers function-space information from a larger teacher to a compact student \cite{hinton2015distilling}. Detector-specific methods distinguish foreground/background regions and localization distributions. Attention transfer and focal/global distillation operate on intermediate representations \cite{zhang2018attentiontransfer,yang2022fgd}, whereas localization distillation treats discrete regression distributions as structured knowledge \cite{zheng2022ld}. Our final training uses response-level class distillation and dense/focused DFL distillation, but no feature imitation or auxiliary detector at inference.

Structural re-parameterization provides the deployment bridge. RepVGG shows that linear branches with batch normalization can be converted into a single equivalent convolution \cite{ding2021repvgg}. We apply the same algebra to a shared/classification-private task decomposition: the stem kernels are concatenated in the channel dimension, the private columns of the box projection are padded with zeros, and the class projections are concatenated. This produces a static deployment graph without gating, dynamic routing, or additional post-processing.
